\documentclass[11pt,a4paper,oneside]{book}
\usepackage{90-faelles/latex/studiearkiv}
\hypersetup{pdftitle={Matematik – Bind 2: Calculus og analyse}}
\begin{document}
\frontmatter
\begin{titlepage}
\centering
\vspace*{30mm}
{\sffamily\color{arkivblaa}\Large MATEMATIK TIL UNIVERSITETSFYSIK\par}
\vspace{16mm}
{\Huge\bfseries Bind 2\\[6mm]Calculus og analyse\par}
\vspace{13mm}
{\Large Formler, begreber og matematisk sammenhæng\par}
\vfill
\begin{minipage}{0.8\textwidth}
Grænseværdier, afledte, integraler, konvergens og flere variable.

Dette bind genbruger de samme LaTeX-kapitler som det samlede matematikkompendium. Planlagte kapitler er markeret som ufærdige. Der indgår ikke separate gennemregnede opgaveeksempler.
\end{minipage}
\vspace{14mm}
{\small Fysik-studiearkiv · Arbejdsudgave\par}
\end{titlepage}
\chapter*{Læsevejledning}
\addcontentsline{toc}{chapter}{Læsevejledning}
Dette bind er en emneopdelt udgave af det samlede matematikkompendium. Den generelle kilde og alle figurer vedligeholdes centralt under \texttt{00-matematik/chapters} og \texttt{00-matematik/figures}. For hele overblikket samt henvisninger mellem bindene anvendes den samlede udgave. Ufuldstændige kapitler er tydeligt markerede.
\begingroup\small\tableofcontents\endgroup
\mainmatter
\chapter{Differentialregning}
\label{ch:diff}
Differentialregning beskriver lokal ændring. Den centrale idé er, at en tilstrækkeligt regulær funktion tæt på et punkt kan erstattes af en affin funktion med en kontrolleret fejl. Hældningen er samtidig en grænseværdi, en lineær afbildning og en lokal følsomhed. Disse tre synsvinkler forbinder de sædvanlige regneregler med fysikkens hastigheder, responskoefficienter og perturbative beskrivelser.

Kapitlet behandler funktioner af én reel variabel. Partielle afledte, Jacobimatricer og differentiabilitet i flere dimensioner hører til kapitel~\ref{ch:flervariabel}. De nævnes kun for at afgrænse, hvilke konklusioner der kan overføres.

\section{Grænseværdier, kontinuitet og lokal adfærd}
\label{sec:graense}
Lad $f$ være defineret i en punkteret omegn af $x_0$. Udsagnet $\lim_{x\to x_0}f(x)=L$ betyder, at der for hvert $\varepsilon>0$ findes et $\delta>0$, sådan at
\[
0<|x-x_0|<\delta\quad\Longrightarrow\quad |f(x)-L|<\varepsilon.
\]
Definitionen kræver ikke, at $f(x_0)$ findes. Kontinuitet i $x_0$ kræver derimod både en funktionsværdi og sammenfaldet $\lim_{x\to x_0}f(x)=f(x_0)$. En venstre- og en højregrænseværdi skal eksistere og være ens, før en endelig tosidet grænseværdi findes. Ved endepunkter arbejder man relativt til definitionsmængden.

Grænseværdier kan føres gennem summer og produkter, når de enkelte endelige grænseværdier findes. For kvotienter skal nævnerens grænseværdi desuden være forskellig fra nul. Ubestemte former som $0/0$ og $\infty-\infty$ er signaler om, at udtrykkets struktur skal undersøges; de er ikke regneregler.

For $h\to0$ betyder $r(h)=o(h)$, at $r(h)/h\to0$. Udsagnet $r(h)=O(h^2)$ betyder, at der findes $C,\delta>0$ med $|r(h)|\le C|h|^2$ for $0<|h|<\delta$. Lille-$o$ udtrykker en forsvindende relativ fejl, mens store-$O$ giver en størrelsesorden med en lokal konstant. Notationen siger ikke noget om en numerisk fejlgrænse, før konstanten er bestemt.

\section{Den afledte og differentialet}
\label{sec:def}
\begin{definition}
Lad $x_0$ være et indre punkt i intervallet $I$. Funktionen $f\colon I\to\R$ er differentiabel i $x_0$, hvis den endelige grænseværdi
\[
 f'(x_0)=\lim_{h\to0}\frac{f(x_0+h)-f(x_0)}{h}
\]
eksisterer. En uendelig hældning tæller ikke som differentiabilitet i denne definition.
\end{definition}
Sekanthældningen måler den gennemsnitlige ændring over et endeligt interval. Den afledte er grænsen af disse hældninger, når intervallet kollapser mod punktet. Den geometriske tangent har ligningen
\[
 y=f(x_0)+f'(x_0)(x-x_0).
\]
Den samme definition kan omformes til
\begin{equation}
 f(x_0+h)=f(x_0)+f'(x_0)h+r(h),\qquad r(h)=o(h).
 \label{eq:linear}
\end{equation}
Denne formulering er særlig vigtig: differentiabilitet er eksistensen af en lokal lineær respons med en fejl, der forsvinder hurtigere end tilvæksten. Differentialet er afbildningen $\dd f_{x_0}\colon h\mapsto f'(x_0)h$. Skrivemåden $\dd f=f'(x)\dd x$ udtrykker denne lineære sammenhæng; den faktiske ændring $\Delta f$ indeholder desuden restleddet.

Hvis $x$ har enhed $[x]$ og $f$ har enhed $[f]$, har $f'$ enhed $[f]/[x]$. Differentialet har samme enhed som $f$. En logaritme eller en eksponentialfunktion skal i en fysisk model have et dimensionsløst argument.

\begin{figure}[htbp]
\centering
\begin{tikzpicture}
\begin{axis}[width=.88\textwidth,height=65mm,axis lines=middle,xlabel={$x$},ylabel={$y$},xmin=-.2,xmax=2.3,ymin=-.8,ymax=4.8,legend pos=north west,legend style={font=\small},samples=100]
\addplot[arkivblaa,thick,domain=0:2.15]{x^2};\addlegendentry{$f(x)=x^2$}
\addplot[arkivrod,thick,domain=0:2.2]{2*x-1};\addlegendentry{Tangent ved $x_0=1$}
\addplot[black,dashed,domain=.4:2.15]{3*x-2};\addlegendentry{Sekant, $h=1$}
\fill[black] (axis cs:1,1) circle[radius=1.5pt];
\fill[black] (axis cs:2,4) circle[radius=1.5pt];
\end{axis}
\end{tikzpicture}
\caption{Sekanten gennem $(1,1)$ og $(2,4)$ har hældning 3; tangenten ved $(1,1)$ har hældning 2. For $h\to0$ går sekantens hældning $2+h$ mod 2. Akserne er dimensionsløse.}
\label{fig:sekant}
\end{figure}


\subsection{Kontinuitet er nødvendig, men ikke tilstrækkelig}
Af \eqref{eq:linear} følger $f(x_0+h)-f(x_0)\to0$, så differentiabilitet medfører kontinuitet. Den omvendte implikation gælder ikke. Funktionen $|x|$ er kontinuert i nul, men har ensidige afledte $-1$ og $1$ og dermed ingen tosidet afledt der.

En afledt behøver heller ikke være kontinuert. Funktionen $f(x)=x^2\sin(1/x)$ for $x\ne0$ og $f(0)=0$ har $f'(0)=0$, mens $f'(x)=2x\sin(1/x)-\cos(1/x)$ for $x\ne0$ ikke har en grænseværdi ved nul. Man må derfor ikke erstatte antagelsen $f\in C^1$ med blot ``$f$ differentiabel'', hvis et argument kræver kontinuitet af $f'$.

\section{Regneregler og sammensatte afhængigheder}
\label{sec:regler}
Nedenstående regler gælder i et punkt, hvor de nævnte funktioner er differentiable. Konstanterne $a,b$ er uafhængige af $x$. For kædereglen skal den indre funktion være differentiabel i $x$, og den ydre funktion i den tilsvarende værdi.

\begin{longtable}{@{}p{0.23\textwidth}p{0.40\textwidth}p{0.29\textwidth}@{}}
\caption{Regneregler. Forudsætninger er en del af formlerne.}\label{tab:regler}\\
\toprule Regel & Formel & Ekstra betingelse\\\midrule\endfirsthead
\toprule Regel & Formel & Ekstra betingelse\\\midrule\endhead
Linearitet & $(af+bg)'=af'+bg'$ & $a,b$ konstante\\[2mm]
Produkt & $(fg)'=f'g+fg'$ & Begge faktorer differentiable\\[2mm]
Kvotient & $\displaystyle (f/g)'=\frac{f'g-fg'}{g^2}$ & $g(x)\ne0$\\[3mm]
Kæderegel & $(f\circ g)'(x)=f'(g(x))g'(x)$ & Sammensætning defineret\\[2mm]
Reciprok & $(1/f)'=-f'/f^2$ & $f(x)\ne0$\\[2mm]
Logaritmisk & $(\ln|f|)'=f'/f$ & $f(x)\ne0$\\[2mm]
Potens af funktion & $(f^\alpha)'=\alpha f^{\alpha-1}f'$ & For reel $\alpha$ er $f>0$ en tilstrækkelig betingelse\\[2mm]
Variabel eksponent & $(f^g)'=f^g(g'\ln f+gf'/f)$ & $f>0$, begge differentiable\\
\bottomrule
\end{longtable}
Produktreglen kommer fra at tilføje og trække $f(x+h)g(x)$ i produktets differens. Efter division med $h$ giver kontinuiteten af $f$ den ene faktor og differentiabiliteten af $g$ den anden. Kædereglen udtrykker, at to lokale lineariseringer sammensættes: den ydre respons skal vurderes ved den indre funktions aktuelle værdi.

Ved tidsafhængighed $q=q(t)$ er $\dd f(q(t))/\dd t=f'(q(t))\dot q(t)$. Prikken betyder differentiation med hensyn til tid, ikke med hensyn til $q$. Denne skelnen forebygger, at en partiel afhængighed forveksles med en total tidsafledt i mekanik og termodynamik.

\subsection{Inverse og implicitte funktioner}
Hvis $f$ er $C^1$ nær $x_0$ og $f'(x_0)\ne0$, er $f$ lokalt invertibel, og
\[
 (f^{-1})'(y_0)=\frac{1}{f'(x_0)},\qquad y_0=f(x_0).
\]
Resultatet er lokalt. En funktion kan godt have flere adskilte monotone grene. Et punkt med $f'=0$ er ikke automatisk et punkt uden en invers, men den inverse kan miste differentiabilitet der.

For en relation $F(x,y)=0$ med $F\in C^1$ og $F_y(x_0,y_0)\ne0$ findes lokalt en funktion $y(x)$, og
\[
 \frac{\dd y}{\dd x}=-\frac{F_x}{F_y}.
\]
Her betegner $F_x,F_y$ partielle afledte, vurderet langs løsningen. Hvis $F_y=0$, er formlen ikke anvendelig; man må undersøge en anden lokal parametrisering eller en mulig singularitet.

\section{Standardafledte med definitionsområder}
\label{sec:standard}
Trigonometriske formler forudsætter radianmål. Ved gradmål kommer en faktor $\pi/180$ ind via kædereglen. De inverse trigonometriske funktioner forstås som deres sædvanlige reelle hovedgrene.

\begin{longtable}{@{}p{0.28\textwidth}p{0.32\textwidth}p{0.30\textwidth}@{}}
\caption{Elementære standardafledte på åbne reelle definitionsområder.}\\
\toprule $f(x)$ & $f'(x)$ & Område eller betingelse\\\midrule\endfirsthead
\toprule $f(x)$ & $f'(x)$ & Område eller betingelse\\\midrule\endhead
$c$ & $0$ & $c$ konstant\\[1mm]
$x^n$ & $nx^{n-1}$ & Heltal $n\ge1$; også nul, hvor udtrykket er defineret\\[1mm]
$x^\alpha$ & $\alpha x^{\alpha-1}$ & $x>0$, $\alpha\in\R$\\[1mm]
$\e^x$ & $\e^x$ & $x\in\R$\\[1mm]
$a^x$ & $a^x\ln a$ & $a>0$ konstant\\[1mm]
$\ln x$ & $1/x$ & $x>0$\\[1mm]
$\ln|x|$ & $1/x$ & $x\ne0$\\[1mm]
$\log_a x$ & $1/(x\ln a)$ & $x>0$, $a>0$, $a\ne1$\\[1mm]
$\sin x$ & $\cos x$ & $x\in\R$\\[1mm]
$\cos x$ & $-\sin x$ & $x\in\R$\\[1mm]
$\tan x$ & $1/\cos^2 x$ & $\cos x\ne0$\\[1mm]
$\cot x$ & $-1/\sin^2 x$ & $\sin x\ne0$\\[1mm]
$\arcsin x$ & $1/\sqrt{1-x^2}$ & $|x|<1$\\[1mm]
$\arccos x$ & $-1/\sqrt{1-x^2}$ & $|x|<1$\\[1mm]
$\arctan x$ & $1/(1+x^2)$ & $x\in\R$\\[1mm]
$\sinh x$ & $\cosh x$ & $x\in\R$\\[1mm]
$\cosh x$ & $\sinh x$ & $x\in\R$\\[1mm]
$\tanh x$ & $1/\cosh^2 x$ & $x\in\R$\\[1mm]
$\operatorname{arsinh}x$ & $1/\sqrt{1+x^2}$ & $x\in\R$\\[1mm]
$\operatorname{arcosh}x$ & $1/\sqrt{x^2-1}$ & $x>1$\\[1mm]
$\operatorname{artanh}x$ & $1/(1-x^2)$ & $|x|<1$\\
\bottomrule
\end{longtable}
En formel gælder kun, hvor den oprindelige funktion er defineret, og de anførte differentiabilitetskrav er opfyldt. Algebraisk forkortelse kan give et udtryk med et større definitionsområde end den oprindelige funktion; det udvider ikke automatisk denne funktion. For $|x|$ er den afledte $\operatorname{sgn}(x)$ for $x\ne0$, men ikke ved nul.

\section{Middelværdisætninger og globale konklusioner}
\label{sec:middel}
\begin{saetning}[Rolle]
Er $f$ kontinuert på $[a,b]$, differentiabel på $(a,b)$ og $f(a)=f(b)$, hvor $a<b$, findes et $c\in(a,b)$ med $f'(c)=0$.
\end{saetning}
Kontinuiteten på det kompakte interval sikrer, at maksimum og minimum antages. Hvis funktionen ikke er konstant, ligger mindst ét af disse ekstrema i det indre. Differentiabilitet tvinger da den afledte til at være nul der. Endepunktskontinuiteten og den indre differentiabilitet udfører altså forskellige opgaver i argumentet.

\begin{saetning}[Lagranges middelværdisætning]
Er $f$ kontinuert på $[a,b]$ og differentiabel på $(a,b)$, findes et $c\in(a,b)$, sådan at
\[
 f(b)-f(a)=f'(c)(b-a).
\]
\end{saetning}
Sætningen følger ved at trække sekantens affine funktion fra $f$ og anvende Rolle. Den siger, at den gennemsnitlige ændring svarer til en øjeblikkelig ændring et sted i intervallet. Den angiver generelt hverken et entydigt $c$ eller en konstruktiv metode til at finde det.

Hvis $|f'(x)|\le M$ på intervallet, følger stabilitetsestimatet
\[
 |f(b)-f(a)|\le M|b-a|.
\]
Det er en Lipschitzgrænse: en begrænset lokal følsomhed giver en kontrolleret endelig ændring. Hvis $f'\ge0$ på et interval, er $f$ ikke-aftagende. Hvis $f'>0$ overalt, er $f$ strengt voksende. Den omvendte slutning er stærkere end nødvendigt: en strengt voksende differentiabel funktion kan have $f'=0$ i enkelte punkter.

Cauchys middelværdisætning gælder for to funktioner med de samme kontinuitets- og differentiabilitetskrav:
\[
 [f(b)-f(a)]g'(c)=[g(b)-g(a)]f'(c)
\]
for et $c\in(a,b)$. Man må først dividere med $g'(c)$ og $g(b)-g(a)$, når de er forskellige fra nul. Denne sætning ligger bag l'Hôpitals regel.

\section{Højere afledte, monotoni og krumning}
\label{sec:hoejere}
Den anden afledte $f''$ beskriver ændringen i hældningen. Klassen $C^n(I)$ består af funktioner med kontinuerte afledte til og med orden $n$ på $I$. $C^\infty$ betyder glat, men betyder ikke nødvendigvis analytisk: en funktion kan være uendeligt differentiabel uden at være lig sin Taylorrække.

\subsection{Ekstrema og stationære punkter}
For en differentiabel funktion er $f'(x_0)=0$ en nødvendig betingelse for et indre lokalt ekstremum. Den er ikke tilstrækkelig. Fortegnsskiftet af $f'$ afgør typen: fra negativ til positiv giver et minimum; fra positiv til negativ giver et maksimum. Intet fortegnsskift giver normalt intet ekstremum, men degenererede tilfælde kræver en direkte undersøgelse.

Hvis $f\in C^2$ nær $x_0$, $f'(x_0)=0$ og $f''(x_0)>0$, er punktet et strengt lokalt minimum; ved $f''(x_0)<0$ et strengt lokalt maksimum. Hvis $f''(x_0)=0$, er testen uafgjort. Mere generelt: er den første ikke-nul afledte efter konstantleddet af orden $m$, og gælder en Taylorudvikling med restled $o((x-x_0)^m)$, giver lige $m$ et ekstremum bestemt af fortegnet; ulige $m$ giver intet ekstremum.

En fuldstændig undersøgelse på et lukket interval skal også omfatte endepunkter og indre punkter, hvor funktionen ikke er differentiabel. En liste over løsninger til $f'=0$ alene er ikke en optimeringsmetode på hele definitionsmængden.

\subsection{Konveksitet og geometrisk krumning}
En differentiabel funktion er konveks på et interval, når dens graf ligger over alle sine tangenter. For $C^2$-funktioner er dette ækvivalent med $f''\ge0$. Konkavitet svarer til $f''\le0$. Et vendepunkt kræver skift mellem konveksitet og konkavitet; ligningen $f''=0$ alene er utilstrækkelig.

For en plan graf $y=f(x)$ med dimensionsløse, lige skalerede koordinater er den geometriske krumnings størrelse
\[
 \kappa(x)=\frac{|f''(x)|}{(1+f'(x)^2)^{3/2}}.
\]
Andenafledte og geometrisk krumning er altså forskellige størrelser. Hvis akserne har forskellige fysiske enheder, skal en fysisk indlejring og skalering fastlægges, før denne geometriske formel kan fortolkes.

\begin{figure}[htbp]
\centering
\begin{tikzpicture}
\begin{axis}[width=.88\textwidth,height=65mm,axis lines=middle,xlabel={$x$},ylabel={$y$},xmin=-2,xmax=2,ymin=-3.5,ymax=3.5,legend pos=north west,legend style={font=\small},samples=150]
\addplot[arkivblaa,thick,domain=-1.9:1.9]{x^3-3*x};\addlegendentry{$f=x^3-3x$}
\addplot[arkivrod,dashed,domain=-1.45:1.45]{3*x^2-3};\addlegendentry{$f'=3x^2-3$}
\fill[black] (axis cs:-1,2) circle[radius=1.5pt];
\fill[black] (axis cs:1,-2) circle[radius=1.5pt];
\fill[black] (axis cs:0,0) circle[radius=1.5pt];
\end{axis}
\end{tikzpicture}
\caption{Afledtens fortegn styrer monotonien. Ved $x=-1$ ligger et maksimum, ved $x=1$ et minimum. Vendepunktet $x=0$ har $f''=0$ med fortegnsskift, men er ikke stationært. Akserne er dimensionsløse.}
\end{figure}


\section{Taylorudvikling og kontrollerede restled}
\label{sec:taylor}
Taylorpolynomiet af grad $n$ omkring $a$ er
\[
 T_n(x)=\sum_{k=0}^n\frac{f^{(k)}(a)}{k!}(x-a)^k.
\]
Det matcher funktionsværdien og de første $n$ afledte i udviklingspunktet. Polynomiet beskriver lokal struktur; det er ikke en global erstatning uden en fejlundersøgelse.

\begin{saetning}[Taylor med Lagrangerest]
Lad $f\in C^{n+1}$ på et åbent interval, der indeholder segmentet mellem $a$ og $x$. Da findes et $\xi$ mellem $a$ og $x$, sådan at
\begin{equation}
 f(x)=T_n(x)+\frac{f^{(n+1)}(\xi)}{(n+1)!}(x-a)^{n+1}.
 \label{eq:rest}
\end{equation}
\end{saetning}
Den valgte $C^{n+1}$-antagelse er en enkel tilstrækkelig betingelse. Hvis $|f^{(n+1)}(t)|\le M$ på segmentet, fås den praktiske fejlgrænse
\[
 |f(x)-T_n(x)|\le\frac{M|x-a|^{n+1}}{(n+1)!}.
\]
Ved blot $f\in C^n$ nær $a$ findes Peanoformen $f(x)=T_n(x)+o((x-a)^n)$ for $x\to a$. Den form giver ikke i sig selv en kendt konstant for fejlens størrelse. Når $f^{(n+1)}$ er kontinuert, er også integralrestleddet gyldigt:
\[
 R_n(x)=\frac{1}{n!}\int_a^x (x-t)^n f^{(n+1)}(t)\dd t.
\]
Integralformen er velegnet til fortegns- og normestimater og gælder også for $x<a$ med integralets sædvanlige orientering.

\subsection{Standardudviklinger omkring nul}
Her er $z$ dimensionsløs, og restleddene gælder for $z\to0$:
\begin{align*}
 \e^z&=1+z+\tfrac12z^2+\tfrac16z^3+O(z^4),\\
 \sin z&=z-\tfrac16z^3+O(z^5),\\
 \cos z&=1-\tfrac12z^2+\tfrac1{24}z^4+O(z^6),\\
 \ln(1+z)&=z-\tfrac12z^2+\tfrac13z^3+O(z^4),\\
 (1+z)^\alpha&=1+\alpha z+\tfrac12\alpha(\alpha-1)z^2+O(z^3).
\end{align*}
Logaritmen og den reelle potens kræver $1+z>0$. De tilhørende uendelige potensrækker har yderligere konvergensbetingelser: logaritme- og den generelle binomialrække har konvergensradius 1, med særskilt analyse på randen; for ikke-negative heltalseksponenter stopper binomialrækken og er et polynomium. Eksponential-, sinus- og cosinusrækkerne konvergerer for alle reelle $z$.

At alle afledte findes, sikrer ikke, at en uendelig Taylorrække repræsenterer funktionen. Funktionen $\e^{-1/x^2}$ for $x\ne0$, udvidet med værdien nul i nul, er glat og har alle afledte lig nul i nul, men er positiv uden for nul. Analyticitet er en ekstra egenskab.

\begin{figure}[htbp]
\centering
\begin{tikzpicture}
\begin{axis}[width=.88\textwidth,height=65mm,axis lines=middle,xlabel={$x$},ylabel={$y$},xmin=-2,xmax=2,ymin=-1.2,ymax=1.2,legend pos=north west,legend style={font=\small},samples=150]
\addplot[arkivblaa,thick,domain=-2:2]{sin(deg(x))};\addlegendentry{$\sin x$}
\addplot[black,dotted,domain=-1.2:1.2]{x};\addlegendentry{$T_1=x$}
\addplot[arkivrod,dashed,thick,domain=-2:2]{x-x^3/6};\addlegendentry{$T_3=x-x^3/6$}
\end{axis}
\end{tikzpicture}
\caption{Taylorpolynomier omkring nul. Den kubiske tilnærmelse følger sinus tættere nær nul end tangenten. For alle reelle $x$ giver Taylor til orden 4, hvor koefficienten til $x^4$ er nul, grænsen $|\sin x-T_3(x)|\le |x|^5/120$. En lille relativ fejl følger ikke automatisk nær sinusfunktionens nulpunkter.}
\end{figure}


\section{Grænseværdier og l'Hôpitals regel}
\label{sec:hopital}
En anvendelig version er følgende: Lad $f,g$ være differentiable på et punkteret ensidigt interval mod et endepunkt $a$, med $g'\ne0$. Antag enten $f,g\to0$ eller $|f|,|g|\to\infty$. Hvis $f'/g'$ har en grænseværdi $L$, endelig eller uendelig, gælder under disse betingelser
\[
 \lim_{x\to a}\frac{f(x)}{g(x)}=L.
\]
Samme princip kan anvendes ved uendelige endepunkter. Tosidige grænseværdier kræver kontrol fra begge sider. Ved gentagen brug skal betingelserne verificeres igen på hvert trin. Hvis kvotienten af de afledte ikke har en grænseværdi, er testen uafgjort; det beviser ikke, at den oprindelige kvotient mangler en grænseværdi.

Reglen differentierer tæller og nævner hver for sig til en ny kvotient; den beregner ikke den afledte af den oprindelige kvotient. Former som $0\cdot\infty$ skal først omskrives til en relevant kvotient. Taylorudvikling er ofte mere informativ, fordi den samtidig viser den første ikke-nul orden og tilnærmelsens fejl.

\section{Fysisk fortolkning og følsomhed}
\label{sec:fysik}
For en tilstrækkeligt glat position $q(t)$ er hastigheden $v=\dot q$ og accelerationen $a=\ddot q$. De har enheder henholdsvis længde pr. tid og længde pr. tid i anden. En grafisk hældning afhænger af aksernes skala, mens den fysiske afledte bestemmes af variablerne og deres enheder.

For en éndimensional konservativ kraft er $F(q)=-U'(q)$. Et ligevægtspunkt opfylder $U'(q_0)=0$. Hvis $U''(q_0)>0$ og massen $m>0$, giver lineariseringen
\[
 m\ddot\eta=-U''(q_0)\eta+O(\eta^2),\qquad \eta=q-q_0,
\]
forudsat tilstrækkelig regularitet, eksempelvis $U\in C^3$ nær $q_0$. Den lineære svingningsfrekvens er $\omega=\sqrt{U''(q_0)/m}$. For $U''(q_0)<0$ er den lineariserede ligevægt ustabil. Ved $U''(q_0)=0$ må højere ordener undersøges. Dette er en generel kobling mellem afledte og stabilitet, ikke en konkret opgaveløsning.

For en målt variabel $x$ giver en lille perturbation $\delta x$ den førsteordensændring $\delta f\approx f'(x)\delta x$. Hvis $x\ne0$ og $f(x)\ne0$, er den relative følsomhed
\[
 \kappa_{\mathrm{rel}}(x)=\left|\frac{x f'(x)}{f(x)}\right|,\qquad
 \left|\frac{\delta f}{f}\right|\approx\kappa_{\mathrm{rel}}\left|\frac{\delta x}{x}\right|.
\]
Nær nulpunkter kan relativ fejl være uegnet, selv når den absolutte fejl er lille. For en stokastisk usikkerhed med lille standardafvigelse $\sigma_x$ giver lineariseringen $\sigma_f\approx|f'(x)|\sigma_x$. Ikke-linearitet, skæve fordelinger og korrelationer kræver en videre analyse i det selvstændige statistikkompendium; formlen er ikke en generel eksakt usikkerhedslov.

\section{Numerisk differentiation og fejlbalance}
\label{sec:numerisk}
Når $f$ kun er tilgængelig gennem evalueringer, kan differenskvotienter approksimere den afledte. For $h\ne0$ er de mest brugte:
\begin{align*}
 D_h^+f(x)&=\frac{f(x+h)-f(x)}{h}
           =f'(x)+\tfrac h2f''(x)+O(h^2),\\
 D_h^0f(x)&=\frac{f(x+h)-f(x-h)}{2h}
           =f'(x)+\tfrac{h^2}{6}f^{(3)}(x)+O(h^4),\\
 D_h^2f(x)&=\frac{f(x+h)-2f(x)+f(x-h)}{h^2}
           =f''(x)+\tfrac{h^2}{12}f^{(4)}(x)+O(h^4).
\end{align*}
Tilstrækkelige lokale regularitetskrav for de viste restled er henholdsvis $C^3$, $C^5$ og $C^6$. De centrerede formler udnytter symmetri til at eliminere ulige eller lige bidrag. Ved rande kræver de funktionsværdier uden for intervallet; brug i stedet en passende ensidig formel eller en fysisk begrundet randudvidelse.

Mindre $h$ er ikke altid bedre. Hvis hver funktionsværdi har absolut fejl højst $\varepsilon_f$, kan fejlbidraget til den centrerede førsteafledte være op til $\varepsilon_f/|h|$. En simpel model for totalfejlen er derfor
\[
 E(h)\lesssim C h^2+\frac{\varepsilon_f}{|h|}.
\]
Her kommer første led fra trunkering og andet led fra afrunding eller datastøj. Balancen giver skalaen $|h|\sim(\varepsilon_f/C)^{1/3}$ op til konstante faktorer, når $C>0$. Dette er et lokalt skalaestimat, ikke et universelt optimalt skridt. Variabelskalering, støjstruktur og ukendte højere afledte påvirker valget.

\begin{figure}[htbp]
\centering
\begin{tikzpicture}
\begin{loglogaxis}[width=.88\textwidth,height=65mm,xlabel={Skridt $h$},ylabel={Fejlmodel},xmin=1e-7,xmax=1e-1,ymin=1e-11,ymax=1e-1,legend pos=north east,legend style={font=\small},samples=100]
\addplot[arkivblaa,dashed,domain=1e-7:1e-1]{x^2};\addlegendentry{$h^2$}
\addplot[arkivrod,dotted,thick,domain=1e-7:1e-1]{1e-12/x};\addlegendentry{$10^{-12}/h$}
\addplot[black,thick,domain=1e-7:1e-1]{x^2+1e-12/x};\addlegendentry{Sum}
\end{loglogaxis}
\end{tikzpicture}
\caption{Illustrativ fejlmodel for centreret differentiation med $C=1$ og $\varepsilon_f=10^{-12}$ i dimensionsløse variable. Minimum ligger ved $h=(\varepsilon_f/2C)^{1/3}\approx7{,}94\cdot10^{-5}$. Kurverne er en analytisk model, ikke målte flydende-kommafejl.}
\end{figure}


Automatisk differentiation beregner afledte ved at føre kædereglen gennem programmets operationer. Den undgår differenskvotientens trunkeringsfejl, men hverken diskontinuerte programgrene, afrunding eller en forkert fysisk model forsvinder af den grund. Symbolsk differentiation, automatisk differentiation og numeriske differenser er forskellige redskaber.

\clearpage
\section{Metodeoverblik og kontrolpunkter}
\label{sec:kontrol}
Begynd med definitionsmængde og ønsket konklusion. Til lokal linearisering bruges differentiabilitet; til en kvantitativ fejlgrænse kræves kontrol over relevante højere afledte. Til global monotoni anvendes fortegnet på den afledte på hele det relevante interval. Til ekstremumskontrol undersøges både stationære punkter, singulariteter og rande.

Ved sammensætninger skal den indre funktions afledte og den ydre funktions evalueringspunkt fremgå. Ved inverse funktioner angives gren og ikke-nul hældning. Ved Taylorapproksimation angives udviklingspunkt, orden, gyldighedsområde og restled. Ved numerisk differentiation afprøves flere skridtstørrelser, så både konvergensområdet og støjgulvet kan identificeres.

En faglig kontrol omfatter desuden:
\begin{itemize}
\item Har alle symboler en definition, og er variable og konstanter skilt ad?
\item Stemmer enhederne, også efter tidsdifferentiation og kædereglen?
\item Er udsagnet lokalt eller globalt, og ligger punktet i det indre eller på randen?
\item Er der divideret med en størrelse, som kan være nul?
\item Er en nødvendig betingelse fejlagtigt brugt som tilstrækkelig?
\item Matcher figurernes akser, kurver og markeringer de beskrevne formler?
\end{itemize}
Disse kontroller er en del af den matematiske argumentation. De erstatter ikke et bevis, men gør skjulte antagelser synlige og forbinder håndbogens formler med forsvarlig anvendelse i fysik.

\chapter{Integralregning og standardintegraler}
\label{ch:integral}
Integralregning beskriver akkumulation, sammenhængen mellem lokale ændringer
og endelige størrelser samt grænseprocesser, hvor mange små bidrag samles.
Et integral kan være en orienteret sum, en lineær funktional eller en
parameterafhængig størrelse. Disse synsvinkler supplerer hinanden, men de
kræver forskellige antagelser. En formeltabel giver derfor kun en del af
svaret: definitionsområde, integrationsgrænser og konvergens skal også
kontrolleres.

Kapitlet behandler først integration af én reel variabel og går derefter
over til uegentlige integraler og parameterafhængighed. Flerdimensionale
koordinatskift og deres Jacobideterminanter hører til hovedemnet om
flervariabel analyse. Ingen af standardformlerne nedenfor er en påstand om
pensum i et bestemt kursus.

\section{Stamfunktioner og analysens hovedsætning}
\label{sec:integral-grundlag}
\subsection{Det bestemte integral som grænseværdi}
For en begrænset funktion $f\colon[a,b]\to\R$, hvor $a<b$, vælges en
inddeling $a=x_0<x_1<\cdots<x_N=b$ og prøvepunkter
$\xi_j\in[x_{j-1},x_j]$. Riemannsummen er
\[
 S_N=\sum_{j=1}^N f(\xi_j)\Delta x_j,\qquad
 \Delta x_j=x_j-x_{j-1}.
\]
Funktionen er Riemann-integrabel, hvis summerne har samme endelige grænse,
uanset inddeling og prøvepunkter, når den største delintervalbredde går
mod nul. Grænsen betegnes $\int_a^b f(x)\dd x$. Kontinuitet på $[a,b]$ er
en tilstrækkelig betingelse; stykkevis kontinuerte, begrænsede funktioner
med endeligt mange brud er også integrable.

Integralet er orienteret: $\int_b^a f=-\int_a^b f$, og
$\int_a^a f=0$. Negativ integrand giver negativt bidrag. Det geometriske
areal mellem graf og akse er derimod $\int_a^b|f(x)|\dd x$. Et nulresultat
for det orienterede integral betyder derfor ikke, at funktionen er nul
eller at arealet er nul. Figur~\ref{fig:integral-fortegn} viser forskellen.
Hvis $f$ har enhed $[f]$ og $x$ enhed $[x]$, har integralet enhed
$[f][x]$; integration med hensyn til tid, længde eller energi er ikke
udskiftelige operationer.

\begin{figure}[htbp]
\centering
\begin{tikzpicture}
\begin{axis}[width=.87\textwidth,height=62mm,axis lines=middle,
 xmin=0,xmax=6.7,ymin=-1.3,ymax=2.3,xlabel={$x$},ylabel={$y$},
 xtick={0,3.14159,6.28318},xticklabels={$0$,$\pi$,$2\pi$},
 legend style={font=\small,at={(.98,.98)},anchor=north east},samples=140]
\addplot[forget plot,draw=none,fill=arkivblaa!18,domain=0:pi] {sin(deg(x))}\closedcycle;
\addplot[forget plot,draw=none,fill=arkivrod!20,domain=pi:2*pi] {sin(deg(x))}\closedcycle;
\addplot[arkivblaa,thick,domain=0:2*pi]{sin(deg(x))};\addlegendentry{$f(x)=\sin x$}
\addplot[arkivrod,dashed,thick,domain=0:2*pi]{1-cos(deg(x))};
\addlegendentry{$F(x)=\int_0^x f=1-\cos x$}
\end{axis}
\end{tikzpicture}
\caption{Orienteret akkumulation: den blå og røde flade under integranden
bidrager med modsat fortegn. Derfor vender $F$ tilbage til nul ved $2\pi$,
selv om det samlede geometriske areal er positivt. Den stiplede kurves
hældning er den blå kurves funktionsværdi. Akserne er dimensionsløse.}
\label{fig:integral-fortegn}
\end{figure}


\begin{samepage}
\subsection{Hovedsætningen og integrationskonstanten}
\begin{saetning}[Analysens hovedsætning]
Er $f$ kontinuert på $[a,b]$, er akkumulationsfunktionen
$F(x)=\int_a^x f(t)\dd t$ differentiabel i det indre med $F'(x)=f(x)$.
Hvis $G$ er en stamfunktion til $f$, gælder
\[
 \int_a^b f(x)\dd x=G(b)-G(a).
\]
\end{saetning}
\end{samepage}
For det første udsagn er differenskvotienten for $F$ gennemsnittet af $f$
over et kort interval. Kontinuiteten gør, at dette gennemsnit går mod
$f(x)$. Det andet udsagn følger, fordi $F-G$ har afledt nul og derfor er
konstant på intervallet. Denne begrundelse forklarer både den lokale
forbindelse til differentiation og forskellen mellem bestemt og ubestemt
integration.

Skrivemåden $\int f(x)\dd x=G(x)+C$ betegner en familie af stamfunktioner,
ikke én entydig funktion. Konstanten $C$ er vilkårlig på hvert sammenhængende
interval, hvor $f$ er defineret. To intervaller adskilt af en singularitet
behøver ikke have samme konstant. En begyndelses- eller randbetingelse kan
fastlægge konstanten, men en algebraisk manipulation alene kan ikke.
En stamfunktion kan eksistere uden en elementær lukket form.

\subsection{Linearitet, sammenligning og integralestimater}
For integrable funktioner på et endeligt interval gælder linearitet og
additivitet, eksempelvis
\[
 \int_a^b(\alpha f+\beta g)\dd x
 =\alpha\int_a^b f\dd x+\beta\int_a^b g\dd x,\qquad
 \int_a^c f+\int_c^b f=\int_a^b f.
\]
Her er $\alpha,\beta$ konstanter, og sidste lighed bruger orienterede
integraler. For $a<b$ medfører $f\le g$, at $\int_a^b f\le\int_a^b g$.
To vigtige størrelsesestimater er
\[
 \left|\int_a^b f\dd x\right|\le\int_a^b|f|\dd x,
 \qquad
 \left|\int_a^b fg\dd x\right|^2
 \le\left(\int_a^b|f|^2\dd x\right)
       \left(\int_a^b|g|^2\dd x\right).
\]
Den sidste ulighed er integralformen af Cauchy--Schwarz for kvadratintegrable
funktioner. Den kontrollerer blandt andet overlap mellem funktioner uden
at beregne integralet eksplicit. Hvis $m\le f\le M$, er
$m(b-a)\le\int_a^b f\le M(b-a)$. For kontinuert $f$ findes et
$c\in[a,b]$ med $\int_a^b f=f(c)(b-a)$; en integrand med spring behøver
ikke antage sin middelværdi.

\section{Substitution, partiel integration og partialbrøker}
\label{sec:integral-metoder}
\subsection{Substitution og orientering}
For kontinuert $f$ og en $C^1$-funktion $g$ på $[\alpha,\beta]$ gælder
\begin{equation}
 \int_\alpha^\beta f(g(t))g'(t)\dd t
 =\int_{g(\alpha)}^{g(\beta)}f(u)\dd u.
 \label{eq:substitution}
\end{equation}
Dette er kædereglen integreret over et interval. Det orienterede udsagn
kræver ikke, at $g$ er monoton: tilbagegående gennemløb annulleres med
fortegn. Hvis man derimod vil omskrive en vilkårlig integralform via en
invers $t=t(u)$, skal en entydig gren vælges, og intervallet må eventuelt
opdeles. Grænser og differentialfaktor skal altid transformeres sammen.

For $u=g(x)$ er differentialet $\dd u=g'(x)\dd x$. Denne notation gør
substitutionen kompakt, men er ikke tilladelse til at kassere faktorer,
dividere med $g'=0$ uden undersøgelse eller blande gamle og nye variable.
Ved uegentlige integraler transformeres først et afskåret, egentligt
integral; grænseprocessen foretages bagefter.

En første genkendelsesregel er at søge en indre funktion, hvis afledte
allerede optræder som faktor. Skift af skala $u=kx+b$ er et koordinatskift,
så faktoren $1/k$ i stamfunktionen er nødvendig, også når kun en tabel
anvendes. Symmetri kan forenkle et integral på $[-a,a]$: en ulige,
integrabel funktion har integral nul, mens en lige funktion giver det
dobbelte af integralet på $[0,a]$.

\subsection{Partiel integration og randled}
For $u,v\in C^1([a,b])$ følger af produktreglen
\begin{equation}
 \int_a^b u(x)v'(x)\dd x
 =[u(x)v(x)]_a^b-\int_a^b u'(x)v(x)\dd x.
 \label{eq:partiel}
\end{equation}
Metoden flytter differentiation fra én faktor til en anden. Et nyttigt
valg lader differentiation forenkle den ene faktor og lader den anden
faktor have en kendt stamfunktion. Dette er en strategi, ikke en garanti
for en kortere beregning. Gentagelse er relevant for polynomier gange
eksponential- eller trigonometriske funktioner.

Randleddet er en del af identiteten. I fysik forsvinder det kun ved
konkrete randbetingelser eller et dokumenteret henfald. På uendelige
intervaller skal både randleddets grænse og integralernes konvergens
undersøges; en formel cancellation af to divergenser er ikke tilstrækkelig.
Ved gentagen brug skal fortegnene og hvilket led der differentieres
holdes adskilt.

\subsection{Rationale funktioner og irreducible faktorer}
For en rational funktion $P(x)/Q(x)$ udføres først polynomiumsdivision,
hvis tællerens grad ikke er mindre end nævnerens. Den egentlige brøk
opløses derefter efter nævnerens reelle faktorer. En gentagen lineær
faktor $(x-r)^m$ kræver alle led $A_j/(x-r)^j$ for $1\le j\le m$.
En irreducibel kvadratisk faktor $q(x)^m$ kræver lineære tællere
$(B_jx+C_j)/q(x)^j$ i hver potens.

Kvadratiske nævnere behandles ved kvadratkomplettering. Positive summer
af kvadrater leder typisk til arctan, mens en forskel af kvadrater leder
til logaritmer og reelle intervalgrene. Partialbrøksopløsningen gælder
på den oprindelige definitionsmængde. Nulpunkter i nævneren må ikke
overses ved forkortelse, og integrationsintervallet må ikke krydse en
pol uden særskilt konvergensanalyse.

\subsection{Trigonometriske og hyperbolske substitutioner}
For $a>0$ kan rodudtryk standardiseres ved
\begin{center}
\begin{tabular}{@{}lll@{}}
\toprule
Rod & Substitution & En anvendelig reel gren\\
\midrule
$\sqrt{a^2-x^2}$ & $x=a\sin t$ & $-\pi/2<t<\pi/2$\\
$\sqrt{x^2+a^2}$ & $x=a\sinh t$ & $t\in\R$\\
$\sqrt{x^2-a^2}$ & $x=a\cosh t$ & $x>a$, $t>0$\\
\bottomrule
\end{tabular}
\end{center}
Den sidste substitution dækker kun den positive ydre gren; den negative
skal behandles særskilt. Identiteten $\sqrt{z^2}=|z|$ forklarer,
hvorfor et kvadratrodstegn ikke uden videre kan fjernes efter et
trigonometrisk koordinatskift.

For rationale udtryk i sinus og cosinus er halvvinkelskiftet
\[
 t=\tan(x/2),\qquad
 \sin x=\frac{2t}{1+t^2},\quad
 \cos x=\frac{1-t^2}{1+t^2},\quad
 \dd x=\frac{2\dd t}{1+t^2}
\]
en systematisk rationalisering. Det kræver delintervaller, hvor
$x\ne(2j+1)\pi$ og den valgte inverse vinkelgren er entydig. Trigonometriske
identiteter, lige/ulige potenser og produkt-til-sum-formler kan ofte give
en enklere metode end dette generelle skift.

\subsection{Reduktionsformler}
En reduktionsformel forbinder en hel familie af integraler. For heltal
$n\ge2$ følger ved partiel integration
\begin{align*}
 \int\sin^n x\dd x
 &=-\frac{\sin^{n-1}x\cos x}{n}
   +\frac{n-1}{n}\int\sin^{n-2}x\dd x,\\
 \int\cos^n x\dd x
 &=\frac{\cos^{n-1}x\sin x}{n}
   +\frac{n-1}{n}\int\cos^{n-2}x\dd x.
\end{align*}
Konstanterne er implicitte i disse relationer mellem stamfunktioner.
For $a>0$ og heltal $n\ge2$ kan
$J_n=\int(x^2+a^2)^{-n}\dd x$ reduceres som
\[
 J_n=\frac{x}{2a^2(n-1)(x^2+a^2)^{n-1}}
       +\frac{2n-3}{2a^2(n-1)}J_{n-1}.
\]
Reduktionen gør det muligt at organisere tabeller efter struktur frem for
at behandle hver potens som en ny metode.

\section{Standardintegraler med gyldighedsbetingelser}
\label{sec:standardintegraler}
Alle tabeller angiver $\int f(x)\dd x=F(x)+C$. Konstanten udelades i
højre kolonne for læsbarhed og vælges uafhængigt på adskilte intervaller.
Medmindre andet er anført, er $a>0$, $k\ne0$, og alle parametre reelle og
uafhængige af $x$. Betingelser vedrører et åbent interval; endepunkter
behandles som grænser. Radianmål anvendes i trigonometriske argumenter.
Logaritmer af fysiske størrelser forstås med en passende referenceskala,
så argumentet er dimensionsløst; en konstant skala ændrer kun $C$.

\subsection{Potenser, logaritmer og eksponentialfunktioner}
\begin{longtable}{@{}p{.23\textwidth}p{.46\textwidth}p{.23\textwidth}@{}}
\caption{Algebraiske og eksponentielle standardstamfunktioner.}\\
\toprule $f(x)$ & $F(x)$ & Betingelse\\\midrule\endfirsthead
\toprule $f(x)$ & $F(x)$ & Betingelse\\\midrule\endhead
$x^p$ & $x^{p+1}/(p+1)$ & $x>0$, $p\ne-1$\\[2mm]
$1/x$ & $\ln|x|$ & $x\ne0$\\[2mm]
$(kx+b)^p$ & $\dfrac{(kx+b)^{p+1}}{k(p+1)}$ & $kx+b>0$, $p\ne-1$\\[3mm]
$1/(kx+b)$ & $\ln|kx+b|/k$ & $kx+b\ne0$\\[2mm]
$\ln|x|$ & $x\ln|x|-x$ & $x\ne0$\\[2mm]
$(\ln|x|)^2$ & $x[(\ln|x|)^2-2\ln|x|+2]$ & $x\ne0$\\[2mm]
$x^p\ln x$ & $\dfrac{x^{p+1}}{p+1}\left(\ln x-\dfrac1{p+1}\right)$ & $x>0$, $p\ne-1$\\[3mm]
$1/(x\ln|x|)$ & $\ln\bigl|\ln|x|\bigr|$ & $x\ne0$, $|x|\ne1$\\[2mm]
$\e^{kx}$ & $\e^{kx}/k$ & $x\in\R$\\[2mm]
$x\e^{kx}$ & $\e^{kx}(kx-1)/k^2$ & $x\in\R$\\[2mm]
$x^2\e^{kx}$ & $\e^{kx}(k^2x^2-2kx+2)/k^3$ & $x\in\R$\\[2mm]
$a^x$ & $a^x/\ln a$ & $a\ne1$\\
\bottomrule
\end{longtable}
For heltallige potenser kan $x^p$ have et større reelt definitionsområde
end den tilstrækkelige betingelse $x>0$ i tabellen. Udvidelsen skal ske ud
fra den konkrete potens, ikke ved at give en generel reel potens en
uvalgt kompleks gren. Tilfældene $p=-1$ og $k=0$ kræver deres særskilte
formler og må ikke indsættes i udtryk med nul i nævneren.

\subsection{Trigonometriske og hyperbolske funktioner}
\begin{longtable}{@{}p{.26\textwidth}p{.43\textwidth}p{.23\textwidth}@{}}
\caption{Periodiske og hyperbolske standardstamfunktioner.}\\
\toprule $f(x)$ & $F(x)$ & Betingelse\\\midrule\endfirsthead
\toprule $f(x)$ & $F(x)$ & Betingelse\\\midrule\endhead
$\sin(kx)$ & $-\cos(kx)/k$ & Alle reelle $x$\\[2mm]
$\cos(kx)$ & $\sin(kx)/k$ & Alle reelle $x$\\[2mm]
$\tan(kx)$ & $-\ln|\cos(kx)|/k$ & $\cos(kx)\ne0$\\[2mm]
$\cot(kx)$ & $\ln|\sin(kx)|/k$ & $\sin(kx)\ne0$\\[2mm]
$1/\cos^2(kx)$ & $\tan(kx)/k$ & $\cos(kx)\ne0$\\[2mm]
$1/\sin^2(kx)$ & $-\cot(kx)/k$ & $\sin(kx)\ne0$\\[2mm]
$\sin(kx)/\cos^2(kx)$ & $1/[k\cos(kx)]$ & $\cos(kx)\ne0$\\[2mm]
$\cos(kx)/\sin^2(kx)$ & $-1/[k\sin(kx)]$ & $\sin(kx)\ne0$\\[2mm]
$\sin^2(kx)$ & $x/2-\sin(2kx)/(4k)$ & Alle reelle $x$\\[2mm]
$\cos^2(kx)$ & $x/2+\sin(2kx)/(4k)$ & Alle reelle $x$\\[2mm]
$\sin(kx)\cos(kx)$ & $\sin^2(kx)/(2k)$ & Alle reelle $x$\\[2mm]
$\sinh(kx)$ & $\cosh(kx)/k$ & Alle reelle $x$\\[2mm]
$\cosh(kx)$ & $\sinh(kx)/k$ & Alle reelle $x$\\[2mm]
$\tanh(kx)$ & $\ln(\cosh(kx))/k$ & Alle reelle $x$\\[2mm]
$\coth(kx)$ & $\ln|\sinh(kx)|/k$ & $x\ne0$\\[2mm]
$1/\cosh^2(kx)$ & $\tanh(kx)/k$ & Alle reelle $x$\\
\bottomrule
\end{longtable}
Polerne i tangens og cotangens er reelle singulariteter, ikke blot en
begrænsning i tabelnotationen. En logaritmisk stamfunktion på hver side
af en pol kan ikke anvendes som én ordinær Newton--Leibniz-evaluering
hen over polen.

\subsection{Rationale udtryk, rødder og inverse funktioner}
\begin{longtable}{@{}p{.27\textwidth}p{.43\textwidth}p{.22\textwidth}@{}}
\caption{Kvadratiske nævnere og rødder; $a>0$.}\\
\toprule $f(x)$ & $F(x)$ & Betingelse\\\midrule\endfirsthead
\toprule $f(x)$ & $F(x)$ & Betingelse\\\midrule\endhead
$1/(x^2+a^2)$ & $\arctan(x/a)/a$ & $x\in\R$\\[2mm]
$1/(a^2-x^2)$ & $\dfrac1{2a}\ln\left|\dfrac{a+x}{a-x}\right|$ & $x\ne\pm a$\\[3mm]
$1/(x^2-a^2)$ & $\dfrac1{2a}\ln\left|\dfrac{x-a}{x+a}\right|$ & $x\ne\pm a$\\[3mm]
$x/(x^2+a^2)$ & $\tfrac12\ln(x^2+a^2)$ & $x\in\R$\\[2mm]
$1/(x^2+a^2)^2$ & $\dfrac{x}{2a^2(x^2+a^2)}+\dfrac{\arctan(x/a)}{2a^3}$ & $x\in\R$\\[3mm]
$1/\sqrt{a^2-x^2}$ & $\arcsin(x/a)$ & $|x|<a$\\[2mm]
$1/\sqrt{x^2+a^2}$ & $\operatorname{arsinh}(x/a)$ & $x\in\R$\\[2mm]
$1/\sqrt{x^2-a^2}$ & $\ln|x+\sqrt{x^2-a^2}|$ & $|x|>a$\\[2mm]
$x/\sqrt{x^2+a^2}$ & $\sqrt{x^2+a^2}$ & $x\in\R$\\[2mm]
$x/\sqrt{a^2-x^2}$ & $-\sqrt{a^2-x^2}$ & $|x|<a$\\[2mm]
$\sqrt{a^2-x^2}$ & $\tfrac12[x\sqrt{a^2-x^2}+a^2\arcsin(x/a)]$ & $|x|<a$\\[2mm]
$\sqrt{x^2+a^2}$ & $\tfrac12[x\sqrt{x^2+a^2}+a^2\operatorname{arsinh}(x/a)]$ & $x\in\R$\\[2mm]
$\sqrt{x^2-a^2}$ & $\tfrac12[x\sqrt{x^2-a^2}-a^2\ln|x+\sqrt{x^2-a^2}|]$ & $|x|>a$\\[2mm]
$1/(x^2\sqrt{x^2+a^2})$ & $-\sqrt{x^2+a^2}/(a^2x)$ & $x\ne0$\\[2mm]
$\arctan x$ & $x\arctan x-\tfrac12\ln(1+x^2)$ & $x\in\R$\\[2mm]
$\arcsin x$ & $x\arcsin x+\sqrt{1-x^2}$ & $|x|<1$\\[2mm]
$\arccos x$ & $x\arccos x-\sqrt{1-x^2}$ & $|x|<1$\\
\bottomrule
\end{longtable}
På den negative ydre gren $x<-a$ er logaritmens absolutværdi afgørende.
En omskrivning til arcosh uden fortegnskontrol kan give den forkerte
afledte. Tabellernes hurtigste lokale kontrol er altid differentiation;
den globale kontrol er fortsat definitionsområde og intervalgrænser.

\subsection{Blandede funktioner og ikke-elementære integraler}
For reelle $k,b$ med $k^2+b^2>0$ gælder
\begin{align*}
 \int\e^{kx}\cos(bx)\dd x
 &=\frac{\e^{kx}}{k^2+b^2}[k\cos(bx)+b\sin(bx)]+C,\\
 \int\e^{kx}\sin(bx)\dd x
 &=\frac{\e^{kx}}{k^2+b^2}[k\sin(bx)-b\cos(bx)]+C.
\end{align*}
Disse formler tillader her også $k=0$, så længe $b\ne0$. For $a>0$
giver kædereglen
\[
 \int x\e^{-ax^2}\dd x=-\frac{\e^{-ax^2}}{2a}+C.
\]
Uden faktoren $x$ findes ikke en elementær stamfunktion. Fejlfunktionen
defineres ved
\[
 \operatorname{erf}z=\frac2{\sqrt\pi}\int_0^z\e^{-t^2}\dd t,
 \qquad
 \int\e^{-ax^2}\dd x=\frac{\sqrt\pi}{2\sqrt a}
             \operatorname{erf}(\sqrt a\,x)+C.
\]
Tilsvarende er en ikke-elementær integralrepræsentation ofte den naturlige
definition af en specialfunktion. Manglende elementær form siger hverken,
at integralet divergerer, eller at det ikke kan beregnes præcist numerisk.

\section{Uegentlige integraler og parameterintegration}
\label{sec:uegentlig-integral}
\subsection{Separate grænser og absolut konvergens}
Et integral med uendelig øvre grænse defineres som
$\int_a^\infty f=\lim_{R\to\infty}\int_a^R f$, hvis grænsen er endelig.
Ved en singularitet i $c\in(a,b)$ skal de to integraler
$\int_a^c f$ og $\int_c^b f$ eksistere hver for sig. På hele den reelle
akse kræves tilsvarende to uafhængige halvakser. En sammenkoblet,
symmetrisk afskæring er en anden grænseproces.

Absolut konvergens betyder $\int|f|<\infty$ og medfører konvergens.
Betinget konvergens betyder konvergens uden absolut konvergens. For
en odd integrand med en simpel pol kan en symmetrisk Cauchy-hovedværdi
være nul, selv om det almindelige uegentlige integral ikke eksisterer.
Hovedværdi skal altid mærkes eksplicit, eksempelvis med
$\operatorname{PV}$; den kan ikke stiltiende erstatte et divergerende integral.

For ikke-negative funktioner kan sammenligning afgøre konvergens.
Modelgrænserne er
\[
 \int_1^\infty x^{-p}\dd x<\infty\ \Longleftrightarrow\ p>1,
 \qquad
 \int_0^1 x^{-p}\dd x<\infty\ \Longleftrightarrow\ p<1.
\]
Ved $p=1$ er divergensen logaritmisk. Et ekstra logaritmisk led kan
ændre tærsklen: $\int_{\e}^\infty[x(\ln x)^p]^{-1}\dd x$ konvergerer
netop for $p>1$. At integranden går mod nul er nødvendigt for mange
regulære positive modelhaler, men generelt ikke en tilstrækkelig test;
figur~\ref{fig:integral-hale} viser to sådanne modelhalers akkumulation.
For vilkårlige kontinuerte, positive funktioner kan selv et konvergent
integral have smalle toppe uden en punktvis grænse mod nul.

\begin{figure}[htbp]
\centering
\begin{tikzpicture}
\begin{semilogxaxis}[width=.87\textwidth,height=61mm,xmin=1,xmax=10000,
 ymin=0,ymax=9.8,xlabel={Afskæring $R$},ylabel={Akkumuleret integral},
 legend style={font=\small,at={(.03,.97)},anchor=north west},samples=130]
\addplot[arkivrod,thick,domain=1:10000]{ln(x)};
\addlegendentry{$\int_1^R x^{-1}\dd x=\ln R$}
\addplot[arkivblaa,thick,domain=1:10000]{1-1/x};
\addlegendentry{$\int_1^R x^{-2}\dd x=1-1/R$}
\addplot[black,dotted,domain=1:10000]{1};
\addlegendentry{Grænseværdi 1 for $p=2$}
\end{semilogxaxis}
\end{tikzpicture}
\caption{To positive modelhaler med integrander, der begge går mod nul.
Integralet af $x^{-2}$ nærmer sig en endelig grænse, mens integralet af
$x^{-1}$ vokser logaritmisk. Den vandrette akse er logaritmisk; tilsyneladende
langsom vækst er ikke konvergens. Størrelserne er dimensionsløse.}
\label{fig:integral-hale}
\end{figure}


\subsection{Oscillation og kontrolleret cancellation}
Dirichlets kriterium er nyttigt ved oscillation: hvis
$G(R)=\int_a^R g(x)\dd x$ er begrænset, og $h$ er monotont aftagende
mod nul, konvergerer $\int_a^\infty h(x)g(x)\dd x$ under sædvanlig
lokal integrabilitet. Dette følger ved partiel integration og kontrol
af randleddet. Resultatet kræver ikke absolut integrabilitet.

Oscillatoriske integraler er derfor følsomme over for hvilken grænse,
rækkefølge og regulering der anvendes. Absolut integrabilitet giver
robusthed over for sådanne omordninger; betinget konvergens giver ikke
samme frihed. En konvergenspåstand bør angive både singulariteter,
haleadfærd og den konkrete definition af integralet.

\subsection{Parameterafledte og bevægelige grænser}
For $f$ og $\partial_\lambda f$ kontinuerte på et passende kompakt
rektangel samt $C^1$-grænser $a(\lambda),b(\lambda)$ er Leibniz' regel
\begin{align}
 \frac{\dd}{\dd\lambda}\int_{a(\lambda)}^{b(\lambda)}f(x,\lambda)\dd x
 &=f(b(\lambda),\lambda)b'(\lambda)
   -f(a(\lambda),\lambda)a'(\lambda)\notag\\
 &\quad+\int_{a(\lambda)}^{b(\lambda)}
                 \partial_\lambda f(x,\lambda)\dd x.
 \label{eq:leibniz-integral}
\end{align}
De første to led beskriver tilførsel eller tab ved intervallets rande;
det sidste beskriver ændring af integranden inde i intervallet. Ved faste
grænser forsvinder kun randleddene, ikke nødvendigvis den indre ændring.

På uendelige intervaller kræves mere end punktvis differentiabilitet.
En praktisk tilstrækkelig betingelse er en integrabel dominator
$g(x)$ med $|\partial_\lambda f(x,\lambda)|\le g(x)$ i en hel omegn af
parameteren, sammen med integrabilitet ved mindst én parameterværdi og
for næsten alle $x$ differentiabilitet i hele parameteromegnen. Domineret konvergens anvendt
på differenskvotienterne giver da differentiation under integraltegnet.
Kontrol ved ét enkelt parameterpunkt alene er ikke denne betingelse.

For grænseovergange $f_n\to f$ næsten overalt er $|f_n|\le g$ med
$\int g<\infty$ en tilsvarende tilstrækkelig betingelse for
$\lim_n\int f_n=\int f$. Uniform konvergens på et endeligt interval er
en enklere tilstrækkelig betingelse. Disse Lebesgue-formulerede resultater
omfatter mange uegentlige integraler, men må ikke erstattes af argumentet
``integranden konvergerer punktvist''. Masse kan flytte mod en rand
eller mod uendelig, uden at integralerne følger den punktvise grænse.

\subsection{Rækkefølge af integraler og parameterfamilier}
Fubinis sætning tillader ombytning af integrationsrækkefølgen, når
$\iint|f|<\infty$. For ikke-negative målbare funktioner giver Tonellis
sætning ombytning også med værdien $+\infty$. Den sidste mulighed er
ikke en påstand om endelig konvergens. Betinget konvergente dobbeltintegraler
kan derimod afhænge af rækkefølgen.

Gaussfamilien
\[
 g_a(x)=\sqrt{a/\pi}\,\e^{-ax^2},\qquad a>0,
\]
har integral 1 på hele den reelle akse. Figur~\ref{fig:integral-gauss}
viser, at parameteren ændrer bredden og højden samtidig. I en kompakt
parameteromegn $0<a_0\le a\le a_1$ kan afledte domineres af et
polynomium gange $\e^{-a_0x^2}$. Der findes ikke automatisk én sådan
lokal dominator, som også kontrollerer hele grænsen $a\to\infty$.
Normalisering, punktvis adfærd og fordeling af masse er forskellige spørgsmål.

\begin{figure}[htbp]
\centering
\begin{tikzpicture}
\begin{axis}[width=.87\textwidth,height=61mm,axis lines=left,
 xmin=-4,xmax=4,ymin=0,ymax=.88,xlabel={$x$},ylabel={$g_a(x)$},
 legend style={font=\small},legend pos=north east,samples=160]
\addplot[arkivblaa,thick,domain=-4:4]{sqrt(.5/pi)*exp(-.5*x^2)};
\addlegendentry{$a=1/2$}
\addplot[black,dashed,thick,domain=-4:4]{sqrt(1/pi)*exp(-x^2)};
\addlegendentry{$a=1$}
\addplot[arkivrod,dotted,very thick,domain=-4:4]{sqrt(2/pi)*exp(-2*x^2)};
\addlegendentry{$a=2$}
\end{axis}
\end{tikzpicture}
\caption{Normaliserede analytiske Gausskerner
$g_a(x)=\sqrt{a/\pi}\,\e^{-ax^2}$. Hver kurve har integral 1 på hele
$\R$, ikke blot i det viste udsnit. Større $a$ giver en højere og smallere
kerne. Figuren illustrerer parameterafhængighed og massens fordeling;
den er ikke et statistisk datasæt. Variable og parametre er dimensionsløse.}
\label{fig:integral-gauss}
\end{figure}


\section{Bestemte standardintegraler og specialfunktioner}
\label{sec:bestemte-standardintegraler}
Nedenstående er familier af generelle identiteter, ikke gennemregnede
opgaveeksempler. Konvergenskravene hører til hver formel. Parametrene
$a,\beta$ er positive, medmindre andet udtrykkeligt angives.
\begin{longtable}{@{}p{.44\textwidth}p{.27\textwidth}p{.21\textwidth}@{}}
\caption{Bestemte integraler med skala-, potens- og oscillationsparametre.}\\
\toprule Integral & Værdi & Betingelse\\\midrule\endfirsthead
\toprule Integral & Værdi & Betingelse\\\midrule\endhead
$\int_{-\infty}^{\infty}\e^{-ax^2}\dd x$ & $\sqrt{\pi/a}$ & $a>0$\\[2mm]
$\int_0^\infty x^{p-1}\e^{-\beta x}\dd x$ & $\Gamma(p)/\beta^p$ & $p>0$\\[2mm]
$\int_0^1t^{p-1}(1-t)^{q-1}\dd t$ & $B(p,q)$ & $p,q>0$\\[2mm]
$\int_0^1x^{p-1}(\ln x)^n\dd x$ & $(-1)^n n!/p^{n+1}$ & $p>0$, $n\ge0$ heltal\\[2mm]
$\int_0^{\pi/2}\sin^m x\cos^n x\dd x$ & $\tfrac12 B(\tfrac{m+1}2,\tfrac{n+1}2)$ & $m,n>-1$\\[2mm]
$\int_0^\infty\e^{-ax}\cos(bx)\dd x$ & $a/(a^2+b^2)$ & $b\in\R$\\[2mm]
$\int_0^\infty\e^{-ax}\sin(bx)\dd x$ & $b/(a^2+b^2)$ & $b\in\R$\\[2mm]
$\int_0^\infty\e^{-ax}\dfrac{\sin(bx)}x\dd x$ & $\arctan(b/a)$ & $b\in\R$\\[3mm]
$\int_0^\infty\dfrac{\sin(bx)}x\dd x$ & $\tfrac\pi2\operatorname{sgn}b$ & $b\ne0$, betinget\\[3mm]
$\int_0^\infty\dfrac{\cos(bx)}{x^2+a^2}\dd x$ & $\dfrac\pi{2a}\e^{-a|b|}$ & $b\in\R$\\[3mm]
$\int_0^\infty\dfrac{\dd x}{1+x^p}$ & $\dfrac\pi p\csc(\pi/p)$ & $p>1$\\
\bottomrule
\end{longtable}
Her defineres gamma- og betafunktionerne ved de angivne integraler;
for $p,q>0$ gælder
\[
 \Gamma(p+1)=p\Gamma(p),\qquad
 B(p,q)=\frac{\Gamma(p)\Gamma(q)}{\Gamma(p+q)}.
\]
Gamma-rekursionen følger af partiel integration, hvor randleddene
forsvinder ved de angivne konvergenskrav. Gaussintegralet kan begrundes
ved at kvadrere integralet og bruge polære koordinater i et ikke-negativt
dobbeltintegral; Jacobifaktoren er afgørende. En fuld videre behandling
af specialfunktionerne hører til kompendiets særskilte hovedemne.

\section{Numerisk integration og dokumenteret fejlkontrol}
\label{sec:integral-numerisk}
Riemannsummer viser definitionen, men er ikke altid den mest effektive
numeriske metode. Figur~\ref{fig:integral-riemann} illustrerer, hvordan
stykvise rektangler approksimerer en glat integrand. Fortegnet af fejlen
kan ved monotoni kontrolleres uden at kende en stamfunktion.

\begin{figure}[htbp]
\centering
\begin{tikzpicture}
\begin{axis}[width=.87\textwidth,height=60mm,axis lines=left,
 xmin=0,xmax=2.15,ymin=0,ymax=1.15,xlabel={$x$},ylabel={$f(x)$},
 samples=120]
\pgfplotsinvokeforeach{0,...,7}{
 \pgfmathsetmacro{\leftx}{#1/4}
 \pgfmathsetmacro{\rightx}{(#1+1)/4}
 \pgfmathsetmacro{\rectheight}{exp(-#1/4)}
 \edef\rectanglecmd{\noexpand\draw[draw=arkivblaa!65,fill=arkivblaa!13]
   (axis cs:\leftx,0) rectangle (axis cs:\rightx,\rectheight);}
 \rectanglecmd
}
\addplot[arkivrod,thick,domain=0:2]{exp(-x)};
\node[anchor=north east,font=\small] at (rel axis cs:.96,.96)
 {$f(x)=\e^{-x}$, $N=8$};
\end{axis}
\end{tikzpicture}
\caption{Venstre Riemannsum for en monoton aftagende, glat funktion på
$[0,2]$. Hvert rektangel ligger over den tilhørende grafdel, så summen er
en øvre approksimation. Højre prøvepunkter ville give en nedre sum.
Inddelingens bredde skal gå mod nul i definitionen af integralet;
figuren viser én endelig inddeling i dimensionsløse variable.}
\label{fig:integral-riemann}
\end{figure}


På et endeligt interval med ensartet skridt $h=(b-a)/N$ er den sammensatte
trapezregel
\[
 T_N=h\left[\frac{f(a)+f(b)}2+\sum_{j=1}^{N-1}f(a+jh)\right].
\]
Hvis $f\in C^2([a,b])$, fås fejlgrænsen
\[
 \left|\int_a^b f-T_N\right|
 \le\frac{b-a}{12}h^2\max_{[a,b]}|f''|.
\]
For lige $N$ og $f\in C^4$ giver den sammensatte Simpsonregel
\begin{align*}
 S_N&=\frac h3\left[f(a)+f(b)
       +4\sum_{\substack{1\le j\le N-1\\j\text{ ulige}}}f(a+jh)
       +2\sum_{\substack{2\le j\le N-2\\j\text{ lige}}}f(a+jh)\right],\\
 \left|\int_a^b f-S_N\right|
 &\le\frac{b-a}{180}h^4\max_{[a,b]}|f^{(4)}|.
\end{align*}
Disse grænser gælder ikke uændret ved singulariteter eller utilstrækkelig
regularitet. Et uendeligt interval kræver både en afskæringsgrænse og en
kvadraturfejl. En ændring af variabel kan isolere singulariteten eller
gøre en hale lettere at kontrollere, men erstatter ikke en begrundelse
for konvergens.

Adaptive metoder vurderer lokal fejl og fordeler evalueringer efter behov.
At to nærliggende skridtstørrelser giver næsten samme resultat er et
nyttigt konsistenstjek, men ikke et generelt bevis: begge kan overse en
smal top eller en hurtig oscillation. Symbolsk beregning, numerisk
kvadratur og analytiske estimater bør derfor supplere hinanden.

\section{Metodevalg, dimensioner og faglig kontrol}
Et hensigtsmæssigt metodevalg begynder med integralets type: bestemt eller
ubestemt, egentligt eller uegentligt, skalar størrelse eller parameterfamilie.
Fastlæg definitionsområde og singulariteter før algebraiske omskrivninger.
Søg derefter symmetri, en afledt indre funktion, en produktstruktur eller
en rational opløsning. Først herefter vælges tabel, specialfunktion eller
numerisk metode.

Kontroller en stamfunktion ved differentiation på hvert relevant interval.
Kontroller et bestemt resultat ved dimensioner, fortegn, størrelsesgrænser
og forventet parameterafhængighed. Ved en substitution skal grænserne,
differentialet og eventuelle grenvalg følges hele vejen. Ved partiel
integration skal randleddet medtages. Ved en grænseovergang eller ombytning
skal en konkret dominans-, konvergens- eller regularitetsbetingelse nævnes.

I fysik repræsenterer integraler blandt andet arbejde langs en valgt
bevægelse, akkumuleret ladning, samlet energi og normalisering af
kontinuerte tilstande. Den matematiske integration fastlægger ikke selv
den fysiske model eller fortegnskonvention. Enhedskontrol og en forklaring
af, hvad der summeres over, er en del af det færdige resultat.

\chapter{Rækker, konvergens og approksimationer}
\planlagt{Definitioner, notation, forudsætninger, centrale resultater og fysikforbindelser udvikles samlet.}
\section{Følger, punktvis og uniform konvergens}
\leavevmode\par
% Redaktionelt: tilføj definitioner, gyldighed, relationer og relevante figurer.

\section{Konvergenskriterier og potensrækker}
\leavevmode\par
% Redaktionelt: tilføj definitioner, gyldighed, relationer og relevante figurer.

\section{Taylor- og Fourierrækker}
\leavevmode\par
% Redaktionelt: tilføj definitioner, gyldighed, relationer og relevante figurer.

\section{Asymptotiske rækker og fejlkontrol}
\leavevmode\par
% Redaktionelt: tilføj definitioner, gyldighed, relationer og relevante figurer.

% Kildekobling: UG-03 i pensumregisteret. Generel konvergensteori her;
% kompleks Taylor- og Laurentteori i emnet Komplekse tal og kompleks analyse. Undgå gentagne teorifremstillinger.

\chapter{Flervariabel analyse}
\label{ch:flervariabel}
\planlagt{Definitioner, notation, forudsætninger, centrale resultater og fysikforbindelser udvikles samlet.}
\section{Differentiabilitet, Jacobimatrix og kæderegel}
\leavevmode\par
% Redaktionelt: tilføj definitioner, gyldighed, relationer og relevante figurer.

\section{Hessian og kritiske punkter}
\leavevmode\par
% Redaktionelt: tilføj definitioner, gyldighed, relationer og relevante figurer.

\section{Implicitte funktioner og Lagrangemultiplikatorer}
\leavevmode\par
% Redaktionelt: tilføj definitioner, gyldighed, relationer og relevante figurer.

\section{Integration i krumme koordinater og Jacobideterminant}
\leavevmode\par
% Redaktionelt: tilføj definitioner, gyldighed, relationer og relevante figurer.

\end{document}
